\section{Bool (Álgebras de Boolean)}

As Álgebras de Boolean são a base da lógica computacional moderna. Estas estruturas estendem os latices distributivos ao introduzir a noção de complemento e elementos extremos (zero e um).

\subsection{Objetos}
A modelação em MATLAB é feita expandindo a estrutura de \textbf{DLat} com a operação de negação lógica e a definição explícita das constantes de topo e base.

\begin{lstlisting}[caption={Operações em Álgebra Booleana}]
B.elements = {0, 1};
B.not = @(a) 1 - a;      
B.and = @(a, b) min(a, b); 
B.or  = @(a, b) max(a, b); 
\end{lstlisting}

\subsection{Morfismos}
Os morfismos nesta categoria devem ser compatíveis com a negação, a conjunção e a disjunção, além de mapear corretamente os elementos neutros do sistema.